% Auto-generated from book/F_hatcher_dictionary.md — do not edit by hand
% Regenerated by scripts/md_to_latex.py

\chapter{Hatcher Parallel Dictionary}
\label{ch:app_f_hatcher}

Condensed from root \texttt{HATCHER\_\allowbreak{}MAP.\allowbreak{}md} for print. Use the free TN PDF: \url{https://pi.math.cornell.edu/~hatcher/TN/TNbook.pdf}

\bigskip
\noindent\rule{\textwidth}{0.4pt}
\bigskip

\section{F.1 Chapter map}
\label{ch:app_f_hatcher:f-1-chapter-map}

\begin{center}
\small
\setlength{\tabcolsep}{4pt}
\begin{tabularx}{\textwidth}{@{}>{\raggedright\arraybackslash}p{0.12\textwidth}>{\raggedright\arraybackslash}X>{\raggedright\arraybackslash}X@{}}
\toprule
Hatcher \emph{Topology of Numbers} & This book (QGA) & Lift \\
\midrule
Ch. 0 Preview & Ch. 0 & Pythagorean \(\rightarrow\) four squares \(\rightarrow\) Hopf glimpse \\
\emph{(none)} & Ch. 1--2 & Quaternions; Hopf fibration \\
Ch. 1 Farey diagram & Ch. 3 & Gauged Hopf lattice; OP1 \\
Ch. 2 Continued fractions & Ch. 3--4 & Phase walks; lattice paths \\
Ch. 3 Symmetries of Farey & Ch. 4 & Left/right multiplications; LFT comparison \\
Ch. 4 Quadratic forms / topographs & Ch. 5 & Flux topographs; OP2 \\
Ch. 5 Classification of forms & Ch. 6 & Types, reduced configs, Magic Islands; OP3 \\
Ch. 6 Representations & Ch. 7 & \(Z\mapsto\) map; OP4 \\
Ch. 7 Class group & Ch. 8 & Composition; class-group analogue; OP6 \\
Ch. 8 Quadratic fields & Ch. 9 & Quaternion algebras; ideal theory \\
\emph{(applications)} & Ch. 10 & Observations; Table T4; OP5 \\
\bottomrule
\end{tabularx}
\end{center}


\bigskip
\noindent\rule{\textwidth}{0.4pt}
\bigskip

\section{F.2 Concept dictionary}
\label{ch:app_f_hatcher:f-2-concept-dictionary}

\begin{center}
\small
\setlength{\tabcolsep}{4pt}
\begin{tabularx}{\textwidth}{@{}>{\raggedright\arraybackslash}X>{\raggedright\arraybackslash}X@{}}
\toprule
Hatcher & QGA \\
\midrule
Farey edge / mediant & Lattice adjacency / candidate mediant (Model) \\
Continued-fraction zigzag & Fiber phase walk; gauge sequence \\
Linear fractional transformation & Unit quaternion left/right multiplication \\
Conway topograph & Flux topograph \\
Separator line & Separator structure / surface \\
Reduced form & Reduced flux configuration \\
Class number & \texttt{class\_\allowbreak{}number\_\allowbreak{}analogue} / island counts (Model) \\
Gauss composition & \texttt{compose\_\allowbreak{}flywheels} (Model; OP6) \\
Ideal class group & Quaternion order ideal class group (Theorem) + Model bridge \\
Representation by forms & Representation by norms + \(Z\mapsto\) flux configs \\
\bottomrule
\end{tabularx}
\end{center}


\bigskip
\noindent\rule{\textwidth}{0.4pt}
\bigskip

\section{F.3 What is deliberately not 1:1}
\label{ch:app_f_hatcher:f-3-what-is-deliberately-not-1-1}

\begin{itemize}
\item No reuse of Hatcher's prose, figures, or exercises.  
\item ``Quaternionic Farey'' is \textbf{defined} for the gauged Hopf lattice (OP1), not assumed unique.  
\item Physics / observation chapters are extra floors, not TN theorems.
\end{itemize}

\bigskip
\noindent\rule{\textwidth}{0.4pt}
\bigskip


